% Generated draft: requires mathematical and source review.
\documentclass[11pt]{article}
\usepackage[margin=25mm]{geometry}
\usepackage{fontspec}
\setmainfont{Latin Modern Roman}
\usepackage{amsmath,amssymb,mathtools,unicode-math}
\setmathfont{Latin Modern Math}
\usepackage{enumitem,needspace,xurl,hyperref}
\setlength{\parindent}{0pt}
\setlength{\parskip}{5pt}
\setlength{\emergencystretch}{3em}
\newcommand{\N}{\mathbb{N}}
\newcommand{\Z}{\mathbb{Z}}
\newcommand{\Q}{\mathbb{Q}}
\newcommand{\R}{\mathbb{R}}
\newcommand{\C}{\mathbb{C}}
\newcommand{\F}{\mathbb{F}}
\newcommand{\A}{\mathbb{A}}
\newcommand{\PP}{\mathbb{P}}
\newcommand{\E}{\mathbb{E}}
\newcommand{\eps}{\varepsilon}
\newcommand{\dd}{\,\mathrm{d}}
\begin{document}
\section*{UnsolvedMath ID: 30002462}
\textbf{Source problem code:} catalogue entry
\subsection{Definitions and mathematical statement}
The supplied record does not define the graph category, the sampling convention for the profile, or whether realizability is exact or asymptotic. The following is therefore a conditional formulation under the finite, simple, undirected graph interpretation with uniformly sampled induced subgraphs; these conventions require confirmation from the source.

Let \(k\geq 1\) be an integer, and let \(\mathcal H_k\) denote the set of isomorphism types of finite simple undirected graphs on \(k\) vertices. For a finite simple undirected graph \(G\) with \(n=|V(G)|\geq k\), define its induced \(k\)-profile by
\[
\pi_k(G)(H)=\frac{\bigl|\{S\subseteq V(G): |S|=k,\ G[S]\cong H\}\bigr|}{\binom{n}{k}},\qquad H\in\mathcal H_k,
\]
where \(G[S]\) is the subgraph induced by \(S\), and \(G[S]\cong H\) means that it belongs to the isomorphism type \(H\). Thus \(\pi_k(G)\) is a probability distribution on \(\mathcal H_k\).

Under these conventions, the first question is to characterize, for each \(k\geq 1\), the set
\[
\mathcal R_k=\left\{p:\mathcal H_k\to[0,1]:\ \sum_{H\in\mathcal H_k}p(H)=1,\ \exists\,G\text{ finite simple undirected},\ |V(G)|\geq k,\ \pi_k(G)=p\right\}.
\]
The second question asks which global properties of \(G\) can be deduced from \(\pi_k(G)\). For a specified graph property \(\mathcal P\) and a realizable profile \(p\), a conclusion that \(G\) has \(\mathcal P\) would require
\[
\forall\,G\text{ finite simple undirected with }|V(G)|\geq k,\qquad \pi_k(G)=p\ \Longrightarrow\ G\in\mathcal P.
\]
The source does not specify the global properties to be considered, whether the order of \(G\) is known, or any intended quantitative conclusions. It also does not distinguish the exact finite realizability question above from realizability by limits of profiles along graph sequences. Those alternatives are unresolved qualifications, and this draft does not identify them with the displayed formulation.
\subsection{Short English statement}
Conditionally interpreting the profile as the distribution of induced subgraphs on a uniformly chosen \(k\)-vertex subset, characterize the distributions realized by finite graphs and the global graph properties they force. The graph conventions, exact versus limiting interpretation, and intended global conclusions need source confirmation.
\subsection{Sources}
\begin{itemize}
\item Supplied local UnsolvedMath registry record, permanent ID \(30002462\), title \emph{Realizability and Consequences of Graph Profiles}, source problem code \texttt{catalogue entry}, in the collection \emph{Oberwolfach Reports Open Problems}. The record gives proposed year 2014 and supplies no individual proposer.
\item Supplied report citation: \emph{Combinatorics (2014)}, Oberwolfach Reports, DOI 10.4171/owr/2014/01, pp.\ 5--90. Original OWR record: \url{https://doi.org/10.4171/owr/2014/01}. The background describes these reports as documenting mathematical workshops and their open-problem discussions. No problem-specific page or passage is supplied.
\item Supplied literature link labelled ``Conlon Fox Sudakov Ramsey bounds'': \url{https://arxiv.org/abs/1002.0045}.
\item Supplied literature link labelled ``Current graph Ramsey survey'': \url{https://www.its.caltech.edu/~dconlon/GraphRamsey.pdf}.
\item Literature assessment dated exactly \texttt{2026-08-21}: this is supplied dated curation, not independently verified. It labels the status open and states that the bounded-degree Ramsey growth conjecture remains open in general, with bounded-chromatic and other special classes known. That assessment concerns a different stated subject from graph-profile realizability; its applicability to this problem is unresolved. The supplied background itself requests independent verification of time-sensitive status and cited claims.
\item The registry labels this problem \texttt{open} and its statement \texttt{corrected\_verified}. Its supplied verification note reports a manual semantic, typographic, and LaTeX review of a meaning-preserving cleanup. These are supplied registry labels and review claims, not independent mathematical or source verification for this draft.
\end{itemize}
\end{document}
